\chapter[Sermon 60. Of the Signs of Antichrist, and the Image of the Beast Raised by Him.]{Of the Signs of Antichrist, and the Image of the Beast Raised by Him.}
\label{sermon:060}
\markright{Sermon 60. Of the Signs of Antichrist, and the Image of the Beast ...}

And he did great wonders, so that he made fire come down from heaven into the earth in the sight of men, and deceived those who dwell on the earth by the means of those signs which he had power to do in the sight of the beast, saying to those who dwell on the earth: that they should make an image unto the beast, which had the wound of a sword, and did live. And he had power to give a spirit unto the image of the beast, and that the image of the beast should speak. And should cause that as many as would not worship the image of the beast should be killed.

\index{Antichrist}He proceeds most diligently to describe Antichrist and his kingdom, which so greatly challenges the faith of Christ and afflicts his church, so that he might be known and avoided by all people.

He says now he shall do great wonders, by which he understands miracles. Some of these are true, and some are false. I call those true miracles which are done in reality and are not by any crafty juggling counterfeited, and which draw me to the truth and set forth the truth. Of this sort, without doubt, were the miracles of the Prophets and Apostles, holy Martyrs, and chiefly of Moses and Christ. These do good to people, do not harm, nor empty people's purses; rather, they glorify God and make the truth to be believed, by drawing people only to God as to the fountain of all goodness. So John testifies of the Lord’s first miracle done in Cana of Galilee, and says: “This is the first sign that Jesus did at Cana in Galilee, and showed his glory, and his disciples believed on him.” This sign was true, and free from any suspicion of trickery; it was a benefit bestowed upon poor people newly married, by which God was glorified. His disciples were moved by this, and the Spirit of God working also inwardly, believed on Christ. So do all true miracles testify Christ to be healthful and beneficial, and therefore alone to be called upon and worshipped. So do John and Peter interpret the signs or miracles that they themselves wrought, as recorded in Acts 3. And of such miracles we have great abundance in the Evangelical, Apostolical, and Ecclesiastical history: and they have no other end than that we must believe in the Son of God, as which alone gives life and all good things.

\index{Papacy}\index{Rome}\index{Zurich}And I call false those signs that are done through devilish craft or enchantment, or by the crafty juggling and subtle trickery of wicked men: such as are the deeds of witches and enchanters; such as were the wise men of Pharaoh, and Simon Magus; and those of which mention is made in Deuteronomy 13. Finally, such were the miracles of the Friars, burned at Bern in Zurich: and it is greatly to be feared that most of the miracles of all Monks and Hermits have been of this sort. Likewise, those are called false also, which although they are done in deed, yet bear witness to a lie against the truth, confirming the Pope to be head of the church, that images are to be worshipped, that we must pray unto Saints, and go on pilgrimages for religion, yea rather for superstition’s sake, that we must worship a new God lying hidden under the form of Bread and Wine, that God must be honored with vows and monastic practices, and if there be any other thing of like sort. With such deceptive signs are filled at this day all temples, churches, and chapels. These have persuaded many even wise men, and do also at this day. Which thing the Lord prophesied should come to pass, saying: there shall arise false Christs, and false Prophets, and shall show great signs and wonders, that even, if it were possible, the elect should be brought into error. And Paul also: the coming of Antichrist, says he, shall be after the working of Satan, with all power, and signs, and lying wonders, and the rest, which is read in 2 Thessalonians 2. And we know that many Bishops of Rome have wrought signs: but that same is not so excellent, but that the Bishop of Rome has confirmed whatever miracles have been wrought in all Christendom, and augmented them with his bulls and indulgences. Doubtless all had a contrary end to the miracles of Christ, and yet have, for they do not profit men, but empty their purses, put men to sundry charges, and lead them away from the faith of Christ, to the faith of Antichrist, confirming his religion, superstition, and doctrine. Neither does he place these his miracles among the last of his arguments, when the simplicity of the gospel is challenged. But if we be wise, we will beware of them, as of a most present pestilence.

\index{Augustine, Saint}And among his miracles, the Lord, through John, recounts that above all things, he causes fire to come down to the earth, and that in his presence. He seems to allude to the story of Elijah, as we read in the fourth book of Kings, chapter one, and as we mentioned in chapter eleven. It was no small miracle, as recorded in the Acts of the Apostles, that the laying on of the apostles’ hands gave the Holy Spirit. Simon Magus also desired the same grace, but he was severely rebuked by Peter the Apostle, as we read in the eighth chapter of the Acts of the Apostles. And it is to be observed (as Augustine has also admonished in the fifteenth book of \textit{De Trinitate}, chapter twenty-six) that the apostles did not give the Holy Spirit, for it is God alone who gives the Holy Spirit. Therefore, at the prayers of the apostles and at the imposition of hands, the Holy Spirit was given from heaven. Wherefore John the Baptist said that he baptized with water, but that Christ should baptize with fire and with the Holy Spirit. And by fire is signified the Holy Spirit.

\index{Primasius}But Antichrist, the Pope, will boast that he, having power given him from heaven, grants the grace of the Holy Spirit. Doubtless, during consecrations, he says that he gives the Holy Spirit. Likewise, in auricular confessions and absolutions, they brag that they give full absolution of sins, which indeed is a great miracle. They lay their hands upon the confessing sinner and say that they absolve him from the punishment and guilt, by the power received from the most holy See of Rome. Primasius, explaining this passage, says: “It is no marvel if that beast, which falsely assumes the name of the Lamb slain and yet living, also fraudulently claims this gift of the Holy Spirit, colorably by imitation, and feigns a donation to his ministers, as we remember that Simon Magus coveted but could not obtain.”

There is another fire also, which Antichrist calls down from heaven, and casts and throws at his enemies, to avenge himself on them: namely, the dart and thunderbolt of cursing. This was terrible to kings, princes, and people. And these have so much feared the thunderbolt of excommunication that they have done and granted many things, which otherwise no one could have obtained from them. The story is known of Emperor Henry IV. For Platina in the life of Gregory VII shows that this emperor was excommunicated by the Pope. After he adds these things about the emperor, he came speedily to Canossa, where the bishop was with Mathilda, and immediately laying aside his royal robes, went barefoot to the gates of the city, and humbly requested to be let in. His entry was denied, which he took in good part, even though the winter was sharp and everything was frozen hard. Remaining three days in the suburbs of the town, and continually asking for pardon, at length, at the request of Mathilda and the Earl of Savoy and the Abbot of Cluny, he was absolved. Frederick Barbarossa, that he might be reconciled to the Pope, laid his neck under his feet to be trodden on, which clearly demonstrates the injuries done to other kings and people by this beast. I will yet tell of one. The Venetians besieged Ferrara, which pays tribute to the Church of Rome, for which cause they were excommunicated by Clement V. Therefore, France Dandalo, who was afterward created Duke, went into France, where the bishop was at that time, to seek pardon for that offense. It was long before he was admitted to come into the Pope’s sight. At last, he was led with an iron chain about his neck to the bishop’s table, like a dog, and there forced to lie under the table among the dogs so long, until the wrath of Clement, without any clemency, had passed, and he removed that ignominy from his country. And he was therefore always afterward called “dog” by his own countrymen, because he had lain like a dog at the Pope’s table to obtain absolution. This is written in Sabellicus in the end of the ninth book, chapter 7. The Pope, when excommunicating, uses tapers of wax burning, which he throws down to the ground from on high, so that even thereby we might perceive that it is he who calls down the fearful fire from heaven upon men on earth. And these things the beast does in the presence of men, namely, with great confidence and boldness, finally to make me afraid and to keep them in awe. For after the same kind of speaking, Paul also commands to rebuke a great man offending before all men, that others may be afraid. 1 Timothy 5.

Nevertheless, the Lord adds the use and effect of Antichrist’s wonders, so that the church may rightly judge them. For he deceives the inhabitants of the earth by reason of miracles, and so forth. He will seduce by these signs or miracles, namely, by that grace of the spirit which he feigns to give, and with those his excommunications by which he would seem to cast men down into hell, those who are more inclined to earthly than heavenly things. And he will lead them from the faith of Christ to his deceptions. Therefore, we must judge the tokens and doctrine of Antichrist, because they seduce people. Cease therefore to marvel at how it has happened that the Pope has allured so many men of wisdom and learning to his side. You have already heard by what means this is done. Therefore, do not always be foolish; learn, take heed, and believe Christ and his Gospel, and cleave to it.

Again he says, that power is given to the beast (truly by God’s just judgment, that according to the apostles’ saying, those who prefer to believe lies rather than the truth may be judged) that he should work those miracles in the sight of the beast. What is it to work miracles in the sight of the beast, but to do them in the presence of all men, boldly and without fear, even to frighten and deceive the beast itself? Here therefore now we hear of two beasts: the beast that does the wonders, and that later beast in whose sight the other, former beast does those signs. Yea, it shall follow hereafter that both the beast and the false prophet, who works these miracles before the beast, by which he also deceives the beast, shall be cast both into everlasting fire. Who therefore is the former and the two-horned beast, but the Pope? The very same is the false prophet also. And who is the beast, in whose sight the Pope works wonders, but the Image of the beast, and therefore a beast also, inasmuch as empire is raised from the beast and governed by the spirit of the beast.

For it follows that the beast sets up an Image of the beast, and that of the same beast which had the wound of a sword and lives: that is to wit, of that old Roman Empire. Now therefore a new Roman Empire is erected, which nevertheless is not plainly called a beast, but the Image of the beast: that is to say, an empire in deed, but which does not come so near the old one by as far as an image differs from the true example. For the old Roman Empire is as it were an example, of which the empire set up by the Roman Antichrist is only an image, representation, shadow, and as it were a dream, having nevertheless some similitude of the same.

I have explained before, and demonstrated through historical accounts, how the old Roman Empire was torn apart and dissolved. In past times, one emperor ruled in the East at Constantinople, and another in the West at Rome or Ravenna. But from the time of Augustulus, for a period of three hundred years or more, there was no emperor of the West. And the lands that had belonged to the emperors were now held by others, and the empire was utterly lost. Therefore, around the year of our Lord eight hundred, when Charles the Great, King of France, came to Rome on Christmas Day, Leo III, Bishop of Rome, placed the crown upon Charles’s head, and the people proclaimed with a loud voice, “To Charles, the Emperor crowned by God, long life and victory!” These events are recorded in all histories, especially in the fourth book of Aventinus’s Chronicles of the Bourbonois.

Again, when this empire seemed to waver and to decline, as if it were about to fall, the Bishop of Rome instituted seven Princes Electors. Some refer this ordinance to Gregory V, who was Pope when Otto was Emperor. And some to Gregory X, who called Ralph of Abpurga to the empire. More will be said about this anon. But the Lord by John says expressly how the beast said to the inhabitants of the earth that they should make an image of the beast. For the Popes have, by speaking, and not by fighting (as appears in the histories of Bishops of Rome, especially of Leo III), erected a new empire. For by preaching, persuading, and practicing, they brought the empire to King Charles. Certainly Platina, in the life of Leo III, says that the Bishop, intending to gratify King Charles, who had deserved well of the church, after solemn service done in the church of Saint Peter, by a loud voice declares Charles Emperor, and crowns him, with the voices and prayers of the people of Rome. \&c.

But now we must examine more diligently wherefore the newly established empire is called by the Pope, the Image of the old beast. And here, indeed, many things could be alleged, but I shall recount only a few. Above all things, it is called the Image, both because it is named the empire itself, and would be taken for the old empire, where in truth it is a name without the thing, and a vain title, lacking that ancient power, majesty, and glory. For unless the emperor possesses the kingdom by his own inheritance, what kingdom shall he have by the name of emperor? Shall he have Rome? Shall he have Italy, the old seat of the empire? Shall he have France, Spain, Hungary, Germany? For although Germany is now considered the seat of the empire, she has her own princes, her own free cities, and which enjoy their privileges, even though they are called imperial. Theodoric of Niem, a German, and a familiar friend of certain popes, who also wrote the lives of certain bishops of Rome, who were last before the council of Constance, in the third book, chapter 43, of his Histories, writes: Of what magnificence the Roman empire is, at least visibly seen in Germany. For you shall have there an archbishop or a bishop, who has of yearly revenues twice as much as the King of Romans receives in all his dominions. And again, a temporal prince, that has more lands than has the emperor. And so forth.

Moreover, in the old Empire, there was a mighty monarch who wielded full authority and was honored by all men as a god on earth. Such as Caius, Domitian, Diocletian, and others. His image represents the Pope, bishop, and king, and is like a certain earthly god, the greatest monarch, with fullness of power. Furthermore, Rome, or the old beast, had a most honorable Senate. So too does the Bishop of Rome have a princely Senate of proud cardinals. For they are, in effect, all princes. The book of the Roman governments recounts the vicar, or lieutenant of the Diocese of Asia (a diocese, in Greek, signifies a disposition, administration, dispensation, government, or jurisdiction), the vicar of the Diocese of Thrace, and of Pontus. Likewise, there was a noble man, president of the governments in Italy, who had many dioceses under him. And no fewer had the lieutenant of France. Just as the Count of Strasbourg, the captain general of the soldiers at Speyer, and the general of the soldiers at Worms acknowledged the Duke of Mainz as a proconsul, so at this day, the bishops of those cities are subjects to the Archbishop of Mainz. The bishops, therefore, seem by the Pope’s ordinance to succeed in the place of the Roman governments. Certainly, you shall see most of these bishops called not only most reverend fathers in Christ, but also most noble and mighty Dukes and Princes of the Empire. And this is also manifest: that the Emperor of the old beast had his legions, the Roman eagles or standards, and most expert and powerful captains of war. But the high bishop and king of Rome has, in that imagery, obedient children, kings and princes in Europe, not to be despised, whom he may command if need require, to stretch forth the secular power. For so thunders Boniface VIII in the first book \textit{De Major}. And obedient, he says, is he who denies the temporal sword to be in the power of Peter, he misunderstands the word of the Lord, saying: “Put up thy sword into thy sheath.” Therefore, both swords are in the power of the church, namely, both the spiritual and the material sword; but this must indeed be exercised for the church, the other by the church. The spiritual by the priest, the material by the hand of kings and soldiers, but at the will and patience of the high priest.

The old beast had his laws written and published daily in a manner new. The Popes, therefore, after the imitation of the imperial laws, have written decretals and often make new laws. Indeed, they say that the voice and precepts or commands of the Pope are to be received and taken as the words of our Lord Jesus Christ, the Son of God, and Apostle Paul. They add moreover that we must stand by the Pope’s determination; that where the Pope is, there is the general council; where the Pope is, there is our common country. He is compelled or reproved by no man, even though he be called an heretic. He has all laws in his breast, or in the scroll of his breast; he may interpret or expound all things. The same ratifies no sentence; and it is in him alone to take away one man’s right and give it to another. He may take away privileges and, at his will and pleasure, not only change bishops, but also depose the emperor himself and declare no sentence of the emperor. All the world is the Pope’s diocese; and the Pope is the ordinary of all, having fullness of power as well in spiritual matters as temporal. For he is Lord of Lords, and has the right of the King of Kings over all subjects. For he has no peer; he is all things, and above all, and it is necessary for salvation to be under the Bishop of Rome. For there is one consistory or judgment seat of God and of the Pope. These things I have taken from their own books, namely, from their Decretals and glosses. There is a book by Antony Russell of Arezzo, on the power of the Pope and the emperor, where you may read innumerable things of the same sort. But of these things which I have noted hitherto, I suppose it is made plain enough, how the Pope, who is here also called the false prophet, has set up the image of the beast.

Hereunto John adds another thing that the empire, once established, and all things set in order, the beast or false prophet moves all that weight, and puts life into the image, so that it can speak: namely, that the false prophet has given it the ability to speak. For unless the pope confirms the election of the King of Romans, he shall not be considered worthy of the name of Emperor. 22. quest. 5. de forma, in the gloss, the emperor swears to the pope, as a client to his lord. You may read the same in the first book, title 9, de jure iurando, in Clementines. Moreover, who does not see how both the emperor and other princes are surrounded by a company of bishops, who inspire them as to what they should speak or do, and how they should behave in all things? For this reason, legates are also sent, called \textit{legati a latere}. And it is not unknown that in all princes’ councils, the spiritual leaders generally have the chief authority. They are, for the most part, chancellors, secretaries, ambassadors, and so on? And the pope and king openly say how he ought to judge all men, but to be judged by no one. Indeed, his officials also usurp this authority for themselves. If there is any assembly, the Bishop of Rome commonly rules by his spirit and governs the chiefest matters, especially matters of religion. For unless the decrees please the fathers, they threaten to abolish whatever the states have decreed. But if a general or national council is called, it is wholly ruled by the pope’s spirit. It speaks and determines as it pleases the pope. For unless it decrees according to the pope’s pleasure, he will seek to abolish everything altogether. For we recently heard that the synod or council is where the pope is. Innocentius the ninth, in the third question, says, “The judge shall be judged neither by the emperor, nor by the whole clergy, nor by kings, nor by the people.” And the gloss on that passage notes that the council cannot judge the pope, and so on. Therefore, if the entire world were to give a sentence in any matter against the pope, it appears that we ought to stand by the pope’s sentence against them all. Indeed, the same glossier in another place says: “The pope, if he will, may dispense against the council, for he is more than the council.” Most truly, therefore, did the Lord say at this present how the beast had power to give a spirit to the beast, and that the image of the beast should speak. For whoever do not show themselves obedient and willing instruments unto this beast in all his affairs are accounted for dead and rotten members, and therefore to be cut off from this vital body. Indeed.

However, in the meantime, lest I should blame any worthy person, or seem to harshly criticize those who have done no wrong, or should be said to fail to acknowledge the goodness of God working in empires, but rather to find fault with it, and to confound and mix together all things both good and evil without discernment, certain things are here by a lengthy, yet necessary digression, to be admonished and better declared. I admonish therefore and repeat, that the Lord our God is the author of empires, and ordains them for the benefit of people; but that the Devil joins himself with God’s good institutions, and according to his evil nature corrupts those good institutions of God, by influencing people’s affections diversely and applying them to evil matters. Whereupon in governments very many things arise which are to be disliked by the godly: such as tyranny, alteration of the state, and similar things. Nevertheless, although God hates all wickedness and cannot allow any evil, we see that he, by his infinite goodness, uses the evil governments of men to the good or profit of his people. For God loves his church exceedingly, and seeks to relieve and comfort all mankind by empires, although not altogether, or in all things commendable.

\index{Daniel (Book of)}I will therefore not deny that sins within the Western Empire were renewed; that is to say, sins the Image of the beast was set up. These seven hundred years, they have many times governed so, that it has easily appeared that God has wrought the health of his people in the governments. Daniel figured by beasts the four Monarchies of the world, which nevertheless did not suppose that all their Princes were beasts, nor did he condemn all Princes, nor did he think that there had been or should be no good thing in them, although the most part were most corrupted. There were found in the old Roman beast (to speak nothing in the meantime of the Princes of Assyria, Babylon, Medes, Persia, or Macedonia) who have set forth profitable laws, set in the books of Justinian. There have been found under that most cruel old beast, who have advanced the true religion of Christ, and defended earnestly the church of God, such as before we said were Constance, Constantine, Theodosius, and diverse others; which come all under the number of the Empire, but not of the beast. But inasmuch as the beast signifies the Empire, so may there be found Princes under the Image of the beast not a few, who have both set forth wholesome laws, and have employed great benefits upon mankind: as have done Charles, Lewis, and Lothaire of Saxony and others. Notwithstanding that they themselves in many things cannot be allowed of the godly. There are found among the later kings of the new Empire, which in power and majesty were not much unlike the old, in virtues not much behind them, but in certain things [illegible]. There are found who have assayed to purge the empire from Popish corruptions, and to bring the Popes under [?Corum]. But with no great or good success. For what the Otthones, Henries, Lodouicks, Fridericks, briefly many French Princes, Saxons, Swedes, Bavarians, and of Austria have been, many notable testimonies of histories do report; which testify that certain Kings both of France and of other realms also, have not bowed their knees to this Baal; or if they have done at any time, yet have they repented, and have showed some token at the least wise, whereby the wise might perceive that they set not much by that beast.

Therefore, all holy and excellent men who have lived throughout time, during the period when the Image of the Beast has reigned, are to be excused. I mean emperors, kings, princes, bishops, states, cities, and the people of the empire and other realms who lived, but were not under the unhappy image of the beast; for they did not offer themselves to the spirit of the beast to be moved and governed by it, nor did they speak explicitly what the beast commanded. Rather, they spoke against the beast and greatly disapproved of his actions. Therefore, just as I have not included the godly princes and good government within the old monarchies, specifically the old Roman Beast, nor condemned them of beastiality [?], so now, when criticizing the Image of the Beast, I do not confuse the good and godly princes and people, and their government which is not evil, with the corrupt actions of Antichrist. For I always except moderate and profitable empires, honest and godly men, however they live under the Image of the Beast, yet not according to the inspiration of the beast or false prophet.

\index{John, Saint}Hereunto I add this also, that the empire was not suddenly established after the will and pleasure of the Bishop, but by diverse spaces of time, sundry attempts, and treasons innumerable: therefore at the length it devolved to an extremity of corruption, and as I may say, bestiality. Whereby it appears that the prophecy of Saint John is to be applied to the things themselves, and to the times, after the state, maliciousness, and corruption of every thing and time.

It is most certain, and by common consent of all historians plainly testified, that in Charles the Great, through the means of Pope Leo the Third, the empire in the west decayed was renewed: and that thus the image of the beast, that is to wit, of the Roman Empire, was erected. And albeit that at this time the empire decayed in the west was restored by the Pope, yet it is evident that the Popes in the beginning of this empire, by certain donations and gifts, did not use so great power as they usurped to themselves afterward, when they had overthrown and deposed certain emperors. For although the donation seemed to be made by King Pippin, and the Pope was then to have received the beginning of his kingdom, yet that he was subject to emperors and kings with the City of Rome also, this same image and other things prove. In the French Chronicles of the Acts of King Charles in the year of our Lord eight hundred and one, it is found written: afterward having set in order the matters of the city and Bishop of Rome, and of all Italy (therefore Italy then also obeyed the Emperor), not only public, but also (mark) ecclesiastical and private (for all the winter the emperor did nothing else), departing from Rome with his son Philip, he came to Spolet. The same author in the Acts of the year eight hundred and sixteen, says he, Stephen, elected in the place of Leo the Third, takes as great journeys as he could to come to the Emperor, sending in the mean time two ambassadors which might treat with the emperor (Louis the Pious) for his consecration. So likewise in the Acts of the year eight hundred and seventeen, it is shown how Paschal being chosen sent an embassy to Emperor Louis. In the Acts of the year 823, the same Bishop stood at the examination and judgment of the emperor. You may find in the Acts of the next year that Emperor Lothair established the matters of Italy and Rome. Yet doth the same author again make mention of the donation of King Pippin, which gave to Saint Peter Ravenna, and Pentapolis, and all the government. Yet doth he make no mention of the donation either of Charlemagne or of Louis the Pious. The 43rd distinct makes mention thereof. I Louis, \&c., in the gloss is written thus: There Louis gives Rome and diverse other things to Saint Peter and to Paschal the Pope. All historians in manner make mention of the donation of the Kings of France. An abridgment of all gathers out of the library Dolaterane in the third book of Geography, in the Acts of Pippin and Charles. Whereby one may easily conjecture what manner of canon is set forth in the 96th distinct, in these words: Constantine the Emperor hath given and granted to the Apostolic See the Crown and all the Imperial dignity in the City of Rome and in Italy, and in the west parts. Which by and by after he discourses with a lengthy exposition out of the life of Saint Sylvester, written (as they say) by Gelasius, in the chapter following. But Antony, Bishop of Florence, denies in his History that this donation remains in any old books. Cusanus and Laurence Valla have impugned the same: neither hath Otto, Bishop of Friesing, in the 3rd chapter of the 4th book of his story, nor Marsilius of Padua in the defense of peace, nor Raphael Volaterane allowed the same, nor many more that I could rehearse. Moreover, in the Chronicles of the Kings of France, set before the story of Paulus Aemilius of the Acts of Kings of France in the year 755, thus you may read: Pippin again entered into Italy, and Aiulf subdued; he gave gifts to Maximus, Bishop of Rome, also the Dukedom of Ravenna of very great lands, lest any man should unthankfully and unjustly take away this largesse from the French Kings, ascribing to the emperor Constantine which Pippin gave to the church of Rome, against the will of the Greek Emperor, affirming the same possessions to be the right of the Empire. From thence Pippin first received and brought into France the ecclesiastical rites of the Romans and ceremonies of songs.

However, the rule of Charles’s empire and his descendants was not very stable or permanent. From the first year that Charles was made Emperor, to the seventh year of Conrad, who was Lewis III’s nephew, by his brother, the last of the house of Charles, there are reckoned about one hundred and nineteen years. For Charlemagne reigned for fourteen years, Lewis for twenty-six, Lotharius for fifteen, Lewis the second for twenty-one, Charles for two years, Charles, surnamed the second, for three, Crassus for twelve, Arnulf for twelve, Lewis the third for ten, and Conrad for seven. Conrad, lying on his deathbed, nominated King Henry, Duke of Saxony, surnamed Falconer. And thus the empire was transferred to the Germans. This Henry, called the first, never went to Italy, nor was he consecrated or crowned by the Pope. His son, Otto, the first of that name, was summoned to Italy and is said to have gone there with a great army, being received at Rome and hailed by the people as Emperor and Augustus. Otto, writing in the sixth book of Histories, chapter seventeen, affirms from the decrees that Pope Leo VIII of that name did consecrate this Otto, the first King of the Germans. For his father, Henry, refused it. Albert Krantz, in the tenth and eleventh chapters of the fourth book of Saxon matters, affirms that Pope Leo made a surrender of all things that the Popes had received from the kings of France and the author defends this surrender as true. However, the keeper of the library testifies that Otto confirmed the donation of the kings of France, Pippin, Charles, and his sons. Moreover, there remains in the decrees a copy of another document, the forty-third distinct, whereby King Otto binds himself to the Pope, promising that he shall not interfere with anything concerning the Pope and the Romans; secondly, that he shall restore all the lands of St. Peter that come into his hands. Let the reader judge what these lands are.

Shortly after this time, around the year of our Lord 996, it is said that by the decree of Pope Gregory V and with the consent of Otto III, the seven prince-electors were assigned, to whom the defense of the church (as says Wimpelingius) and the Roman Empire was committed. All historians and writers agree on this, including the Italians Blondus, Platina, Sabellicus, Volateranus, Egnatius and others; Germans, Albertus, Nauclerus, Carrion, Functius, and several others. However, some have made no mention of this ordinance. Therefore, Aventinus in the fifth book of Chronicles, folio 510.707, says that he knows [?I cannot tell how certainly] that after the death of Frederick II, the electors were instituted and confirmed by Gregory X. But however that matter stands, it is certain that there have been many among the seven prince-electors who were fervent and earnest in true religion, and excellent in all kinds of virtues, especially the seculars, as they term them, who have greatly disliked the tyranny and impiety of the Popes of Rome, so much so that they have stoutly and often withstood them. Our age is undoubtedly indebted to this order or state, because a good part of the preaching of the holy gospel has been reformed, which both they and other princes of Germany, most worthy of praise, valiantly defend and maintain [?by God’s inspiration]. The Lord increase in them, and in other godly princes throughout the world, his gifts, and mercifully keep and preserve them. But to return to the progress and order of the history, it is certain that immediately after Gregory V, the Devil invaded the see of Rome. Neither could Platina dissemble this thing, a writer of Popes’ lives known to all men, who has very favorably spared his lords and masters, and many times has covered their abominable acts; yet, writing of the successor of Gregory V, Sylvester II, before called Gilbert, a monk of Florey, he says that he deserted his monastery and followed the Devil, giving himself wholly to him. And immediately he adds: Gilbert, moved by ambition and a devilish desire to rule, through bribery first obtained the archbishopric of Reims, then of Ravenna, and afterward, with greater entreaties, the Devil furthering him, he obtained to be pope; yet under this condition that after his death he should be the Devil’s wholly. \&c. He who would know the full story, and abridgment taken from Antoninus, Nauclerus and others, let him read the ninth book of Funccius’ Chronicles, under the year 998. Benno, a Cardinal, supposes at this time to be fulfilled those thousand years, after which the Devil, breaking loose, began again to rage in the world, of which certain things will follow in the twentieth chapter of this book. Nevertheless, I shall not refuse to gather here certain things from this Benno Cardinal, and briefly to recite them here for the declaration of our matter.

Therefore, in the life and acts of Hildebrand, called Gregory the Seventh, one Gerbertus, who had infected the city with sorcery, (he says) after the thousand years fulfilled, coming up out of the pit of God’s permission, was Pope for four years, and changing his name, was called Sylvester the Second. And after Gilbert, for twenty-five years (I suppose it should read thirty-two years), and how they reigned, histories testify, and that very evil Theophilactus’s scholar violently achieved the seat, called Benedict the Ninth. He had a dear friend and privy to all his doings, one Gratiane, Archpriest of St. John, near the Latin gate. To whom Hildebrand, a monk of Cluny, forsaking his abbey, did familiarly cleave and became a familiar friend of his. But Benedict, fearing himself, sold his seat to Gratiane, the master of Hildebrand, receiving from him five hundred thousand pounds, which promoted to the office was called Gregory the Sixth. Nevertheless, they shortly had a third Pope, Sabinus, and he was called Sylvester the Third.

\index{Repentance}The emperor, therefore, Henry the Second, a godly man, valiant, wise, and stout, going to Rome to purge the church (for as yet the bishops used not full authority), compelled Benedict or Theophilacte the Magician to flee, cast Gregory in prison, and sent Sylvester away to his old bishopric. And he holding a council, placed the bishop of Bamberg, whom he called Clement, in the seat, of whom also he received the crown. And he brought Gregory with his disciple Hildebrand with him into Germany. In the meantime, Benedict returning to Rome from flight, vexes Clement, and with much enchantment infects the city; and by letters received from Hildebrand out of Germany, he learns what is done in the emperor’s court. Gregory dies there in prison, and left Hildebrand his heir both of his false packing and of his money. Clement dies also, whom Damasus the Second succeeds immediately, but straightways poisoned; by reason of the tumult that was in the city, the emperor sends Bruno (bishop of Toul, come of the noble house of the earls of Holstein), a worthy man. Here Beno annexes: in whose train, through the overmuch sufferance of the emperor, Hildebrand was permitted to return; by this permission to subvert both the bishopric and empire under pretence of religion. And this Beno herein was a true prophet, which says also in the story of Hildebrand: and telling Bruno many things, by the way crept into his favor; and as soon as he came to Rome, obtained of him that he was made one of the keepers of St. Peter’s altar. And in a short time he filled his coffers. And he also reconciled his old lord and master Benedict, feigning repentance deceitfully to Leo the Ninth (for so Bruno being made Pope was called), and through the counsel of Benedict, otherwise called Theophilacte, he armed Leo against the Normans and betrayed him unto them. The Germans therefore slain by treason, scarcely the Pope escaped all desolate. This says Beno. And certainly it is that this monk Hildebrand, from that time forward aspired to get the seat; and in the meantime, whilst it was governed by others, he incensed and ruled the Popes, as Leo the Ninth, Victor the Second, Stephen the Ninth, Benedict the Tenth, Nicholas the Second, and Alexander the Second. But they smell of Hildebrand’s style, that are set forth in the name of Leo, Nicholas, and Alexander. But at length, he himself climbed into the chair, in which he so used himself that no man unless he were stark blind but might see that his devilish government had requited most abundantly Henry the Fourth, the son of Henry the Third, his father’s carrying of him into Germany. And he began openly and impudently to take upon him the power of the emperor. Neither can it be told in few words, in what detestable wise this beast did afflict both the emperor and empire, all the while he was Pope, for the space of twelve whole years. An abridgement of that story has John Functius compiled in the tenth book, under the year of our Lord 1074.

I know that Platina, and many Italian writers, yes, and some Germans also, do highly commend the religion and virtues of this Gregory the 7th. By this, the Papal tyranny, under the guise of religion, is wonderfully augmented and confirmed, and many are blinded. Yet, through God’s grace, it has come to pass that men of grave authority, religion, and virtues have fairly and well stripped off the disguise from this beast. Therefore, Synods and Councils are not to be condemned; first, the Council of Mainz, where 19 famous Bishops were present. Then, a Synod was assembled at Brixia, with 30 Bishops and most of the nobles of Germany and Italy. There was also a Council assembled at Worms, where King Henry was present, and all the German Bishops (except those of Saxony) deposed the Pope from his office. The Epistles and fragments of these Councils are found in the Chronicles of Verspergen, chiefly. He is openly accused by these of all wickedness and ungodliness, of hypocrisy and cruelty. We have rehearsed a little before what Cardinal Benno, a writer of his time, has committed to writing. There remain also testimonies of Sigisbert, an old writer, concerning this Pope. Anyone who wishes may read the 5th book of Aventinus, from page 162 and onward, and also the preface of the 6th book. The same author, in the 7th book, reporting the words of Eberhard, Bishop of Salisbury, who was present at the Council of Regensberg, says that Hildebrand, 170 years under the guise of religion, laid first the foundation of the Antichrist’s kingdom. This wicked war he himself first began, which his successors have continued to this day. First, they excluded the Emperor from the Pope’s election and referred it to the people and priests of Rome. Afterward, they mocked and thrust out [?]; they go about now also to bring us into subjection and bondage, so that they might reign alone. And the things that follow declare that there have not lived many Popes more bold and impudent than this, who have advanced more highly the majesty of the seat. He excommunicated Emperor Henry the 4th and deprived him of the imperial dignity; moreover, he stirred up his subjects against him and absolved the rebels and traitors from their oaths of fidelity. And he himself, like a monarch, gave the Crown of the Empire to others at his pleasure. The power and treasure of the Empire have been so worn and wasted, both by civil and foreign wars, that for many years now the kings of Germany have neither been able to recover their strength nor yet to resist the most arrogant tyranny of the Popes. And thus, at last, the Pope has become a monarch, and Emperors, Kings, and Princes are made their clients and wards.

When Gregory VII was dead, Mokes of Hildebrand’s sect and faction succeeded him—of his manners and corrupted nature, as if heirs and sons who remain entirely within his character—Victor II, Urban II, Paschal II, and Gelasius II. Paschal II caused the son, Henry V (oh, wicked and detestable parricide), to war against the father, the miserable Henry IV. Shortly after, Gelasius II and Callistus II excommunicated Henry V as well, and did not cease to harass this prince until they had wrested from him the gift or election of the Bishopric. And that to the great and inestimable profit of the See of Rome, and to the unrecoverable loss of Germany, and so forth. These things are described more fully by Urspengen [?] in the Chronicle of the year 1122.

In the times following, the audacity and power of the Popes has increased hourly, and the German kings have resisted them stoutly enough, but yet with small success. Where in the mean season we must remember the words of the Lord, uttered by Daniel, saying: and there shall arise a King of a shameless face, and understanding propositions, and his strength shall be fortified, but not in his own force: and it cannot be believed how he will destroy all things, and he shall prosper, and do \&c.

I will therefore touch briefly what things in the times following the Popes have attempted against kings, and boldly done for the establishing of their empire and monarchy. Pope Alexander III excommunicated Frederick I, called Barbarossa, and trampled him under his feet. And when the good prince said how he showed this obedience to St. Peter, the beast exclaimed, setting himself also before Peter, and said, both to me and to Peter, and stamped on him. Pope Innocentius III could not abide, much less allow Philip, the son of Frederick, created Emperor; but commanded the electors to choose another, I mean Otto, Duke of Saxony, whom notwithstanding shortly after he excommunicated also. That proud beast said that he would take from Philip the imperial crown, or lose his Apostolic mitre. Unto Innocent are ascribed those most proud words, which are read in the decretal of Gregory IX, de Elect., in title 34, chapter 6. on this wise: that the princes have right and authority to choose a king, and afterward to advance him to be Emperor, we acknowledge, as we ought, as to whom of right and ancient custom it is known to appertain; especially since that such right and authority came unto them from the Apostolic Seat, which translated the Roman Empire from the Greeks to the Germans in the person of great Charles. However, the princes must know again that the right and authority to examine the person chosen king, and to be promoted to the empire, belongs unto us, which do anoint, consecrate and crown him, \&c. The same is written in the first book, title 33, de maior. and obedient., writing to the Emperor Constantine. So much diversity, says he, as there is between the sun and the moon, so great a difference is there between Popes and Kings, in God’s name.

But Emperor Frederick II, nephew to Barbarossa, an excellent prince, was excommunicated by many Popes: Honorius III, Gregory IX, and Innocentius IV. Indeed, Gregory IX, while Frederick, that excellent prince, waged war in Syria with the Sultan, invaded and occupied Frederick’s provinces. There were most cruel wars and discords between the Popes and this Frederick. The same Innocentius IV excommunicated Conrad IV of that name, son of Frederick II, and stirred up the Prince of Thuringe against him. And when Emperor Conrad was dead, the Pope obtained the goodwill of the Neapolitans to yield themselves to the See of Rome. Conrad had left a son and heir, Conradine, and Manfred his bastard brother, who would be called king of Sicily. Therefore, Pope Urban IV (some say Clement IV) against Manfred summoned Charles, brother to King Louis the French king, Earl of Provence and of Saintonge, to come with an army into Italy, and called him King of both Sicilies. He overcame and slew Manfred at Benevento and received the kingdoms of Sicily from the Pope to do him homage. But Conradine, Duke of Swabia, accompanied by Frederick, Duke of Austria, led out of Germany a well-equipped army into Italy against Charles for the recovery of his old and paternal kingdom. But vanquished by Charles at Lake Fucino, he was taken with Duke Frederick. It is said that 12,000 were slain. The occasion of so great an evil were the Popes, chiefly Clement IV, who, when questioned by Charles, the worthy prince, about what he should do with his prisoners, answered in a way that the Frenchmen understood meant they must suffer. Therefore, he put them both to the sword. In this, the house and posterity of the most noble Dukes of Austria and Swabia is said to have failed. Paulus Aemilius discusses this matter more at length in the 7th book of French Acts, and Aventinus in the 7th book. But even so, the ire and fury of those most holy fathers could not be pacified, so that the most noble Dukes of Swabia had, for God’s glory and the common good, most godly and most constantly resisted the Roman bishops, wolves, I would have said.

But these parricides and bloody wars displeased all good people everywhere, and especially wise and godly princes, so that they understood how they must avoid that empire and flee from it as from the plague. For it was not only a shadow, but would also utterly consume the yearly revenues and treasure of whoever received the office. Now it was known throughout the world what the most valiant and excellent princes of Germany had endured for 119 years, from Henry IV to the sons of Frederick II, because of the bold ambition and incredible malice of the Popes; and that many of them had lost both their lives and their ancient kingdoms, and most excellent liberty.

And there was the empire without any emperor for several years, which I am accustomed to call a desolation of the kingdom or empire. For the popes with their invincible and intolerable pride and tyranny had so weakened the force of the emperors, that the empire seemed subverted and destroyed; neither could there easily be found anyone who would take it up or thought it worthy to be desired. At last, at the commandment of Gregory X, who held a council at Lyons, Count Rudolph of Habsburg was chosen. Although he did not refuse the offered office, yet when often requested to come to Rome, he is said to have answered: “The wayward steps of feet do fear me sore,” meaning by this saying that he did not trust the popes, who by their schemes had destroyed both many princes of Germany and also innumerable people coming to Rome. And this Rudolph is said to have been crowned king in the year of our Lord 1273, the 200th year after Gregory VII. And so long a time lasted the conflict of popes and emperors. A little while after, while Albert, the son of Rudolph, was chosen emperor, and the election was referred to Boniface VIII of that name, he stoutly rejected it, and showed immediately in word and deed that he was both pope and emperor, who rightfully held both powers. Which I explained in sermon 58, and the same Albert Krantz declares exceedingly well in book 8, chapter 36, of Saxon matters. In the place of King Albert, was substituted Henry, prince of Luxembourg. But what authority over him and the empire challenged Clement V, he who wishes, may know from the Clementines. For there is a long treatise thereof in book 2, title 9. I could also rehearse many other similar things of Pope John XXII and others, if I did not think it superfluous.

For these things which I have discussed up to this point, it appears sufficiently clear that the popes themselves, with a mischievous boldness, have seized imperial power, boast of being monarchs, and abuse the service and ministry of kings, treating them as wards and clients, yet pretending the name of sons, so that they may have them the more obedient. For so Gregory VII wrote to Géza, king of Hungary, as can be read in chapter 17 of this book, sermon 75. Yet understand in the meantime that the great part of princes and nobles have not known this beast, but have rather opposed him, and therefore have not come into the number of the beast. But inasmuch as they lived under the Empire, they were far estranged from the beast.

By this I would have them addressed, who will exclaim and say: who can take it well to have the holy Empire called the Image of the beast, and so many noble Kings and Princes, Cities and worthy people? But I neither ought nor will change the manner of speaking which the scripture uses. These are the Lord’s words all, which Daniel in old time, and now Ibon, have revealed to us; but I may except and excuse those who are excused by the testimony of scripture. The way is ready and brief: whosoever will be free from the beast, let him take heed that he be not influenced by the Pope’s spirit; and that he speak not and do not that which the Pope commands against godliness. Let him rather be ruled by the spirit of Christ; and so it shall come to pass that, dwelling in the midst of Babylon, he shall not live after the iniquities of Babylon, but in the Kingdom of Christ.

It follows: and the beast shall cause that whosoever shall not worship the Image of the beast, shall be slain. And it is all one offense to worship that old beast and to worship the Image of the new beast. Of the worshipping of him, I have spoken a little before. Therefore they worship the Image of the beast, which admit the decrees and those ordinances of the seat and Empire, speaking the inspiration of the beast: which allow the Roman religion, which fall to the kissing of the feet, and show themselves in all things obedient children of the seat, and are faithful to the popish Empire. Now if any will not be such a one, and would be content with Christianity, would abhor Rome the seat of the beast, and detest the Image of the beast, he like a church robber and traitor is judged unworthy of life. There is a canon in the fifth book of Decretals, the seventh title of heretics, wherein without any circumstance of words, Lucius the third of that name, determines plainly that heretics are struck with an everlasting curse, whosoever believe and teach otherwise of the Sacraments than the church of Rome believes and teaches. He commands moreover that such being deprived of all dignity, shall be committed to the judgment of the secular powers to be punished with due correction. But if the temporal magistrate will not punish and so defend the church, then he shall be also deprived of all honor, etc. But why do I tarry in rehearsing these things? All men at this day know and see what things are done daily. They are condemned, exiled, excommunicated, shut up in prisons, vexed with sundry torments, at the length also cruelly slain, whosoever shall refuse to worship both the beast and his Image. The Lord Jesus, the true King and Bishop of his church succor us, and restrain the cruelty of the ungracious beast. Amen.
